\begin{recipe}{Bloemkoolquiche}
\kicker[Hoofdgerecht,Vlees,Hartig]{Hoofdgerecht · vlees · hartig}
\index[register]{Bloemkoolquiche}
\index[register]{Bloemkool}
\index[register]{Spek}

\meta{1 uur 40 \dvd Voor 6 porties \dvd Oven 180\,°C}

\lettrine{D}{eze} quiche maak ik al jaren regelmatig voor de familie. Stevig
deeg, een vulling van bloemkool en spek, en veel kaas. Wie het snel wil
doen gebruikt bladerdeeg in plaats van zelfgemaakt deeg.

\blockrule{Ingrediënten}

\noindent\textit{Deeg}
\begin{ingredients}
  \ing{300 g}{bloem}
  \ing{150 g}{roomboter}
  \ing{2}{eieren}
  \ing{1 tl}{zout}
\end{ingredients}

\medskip
\noindent\textit{Vulling}
\begin{ingredients}
  \ing{1}{bloemkool}
  \ing{250 g}{gerookte spekreepjes}
  \ing{1}{ui}
  \ing{200 g}{crème fraîche}
  \ing{4}{eieren}
  \ing{150 g}{oude kaas, geraspt}
  \ingb{versgemalen zwarte peper}
\end{ingredients}

\blockrule{Bereiding}
\tip{Bedek de zijkanten van het deeg tijdens het voorbakken met aluminiumfolie
zodat ze niet naar beneden zakken.}
\begin{steps}
  \item Kneed de deegingrediënten een paar minuten door. Maak een bol, wikkel
        in folie en laat 15 minuten rusten in de koelkast.
  \item Verdeel de bloemkool in roosjes en kook ze 3 minuten in ruim kokend
        water. Laat uitstomen op een bord.
  \item Snipper de ui. Bak de spekreepjes in een koekenpan tot ze beginnen te
        kleuren, voeg de ui toe en bak tot de ui glazig is.
  \item Voeg de bloemkool toe en bak kort mee. Breng op smaak met peper en roer
        goed door zodat de smaak van ui en spek over de bloemkool trekt. Haal
        van het vuur.
  \item Verwarm de oven voor op 180\,°C (boven- en onderwarmte).
  \item Rol het deeg uit tot de grootte van de springvorm plus de rand. Vet de
        springvorm in met boter en bestrooi licht met bloem. Leg het deeg in de
        vorm en laat het een klein beetje over de rand uitsteken. Prik met een
        vork wat gaatjes in de bodem: dat voorkomt dat het deeg tijdens het
        voorbakken opbolt.
  \item Voorbakken: bak de bodem 15 minuten.
  \item Haal de vorm uit de oven. Strooi een laagje kaas op de bodem, schep de
        bloemkoolmengsel erin en strooi tussendoor ook wat kaas. Eindig met een
        laagje kaas bovenop.
  \item Klop de 4 eieren met de crème fraîche tot een glad mengsel en giet dit
        over de vulling.
  \item Bak de quiche 45 minuten op 180\,°C tot de vulling gestold is.
  \item Laat de quiche 10 minuten rusten voor het aansnijden.
\end{steps}

\heroimagefade{images/quiche_bloemkool}
\end{recipe}
